\section{Recursive Bisection Method}

Recursive graph bisection, originally developed for distributing computational workloads across parallel processors~\cite{karypis1998metis}, adapts naturally to redistricting. The method represents a state's geography as a mathematical graph: census tracts become nodes, geographic adjacency defines edges. An algorithm then partitions this graph into $k$ districts of equal population by recursively splitting it in two.

\subsection{Algorithm Overview}

The process begins with 85,000 census tracts (2020 data) forming a single connected graph. To create $k$ districts, the algorithm performs $\lceil \log_2 k \rceil$ rounds of bisection:

\begin{enumerate}
\item \textbf{Graph construction}: Census tracts become nodes weighted by population; edges connect adjacent tracts based on shared boundaries (queen contiguity).
\item \textbf{Edge weighting}: Each edge $(i,j)$ connecting adjacent tracts receives weight:
\[
w(i,j) = \frac{\alpha \cdot s_{demo}(i,j)}{d(i,j) + \epsilon}
\]
where $d(i,j)$ is centroid distance (km), $s_{demo}(i,j) \in [0,1]$ measures demographic similarity via cosine similarity of racial/ethnic composition vectors, $\alpha$ is a scaling parameter balancing VRA objectives with compactness, and $\epsilon = 0.1$ prevents division by zero for adjacent tracts. Higher weights between demographically similar tracts preserve communities of interest, addressing Voting Rights Act requirements without explicit racial targeting. Full specification and weight distribution statistics provided in Supplementary Materials.
\item \textbf{Bisection}: The METIS graph partitioning library~\cite{karypis1998metis} splits the graph into two balanced parts, minimizing the total weight of cut edges. This optimization produces compact districts by preferring cuts that separate dissimilar communities over those that split cohesive neighborhoods.
\item \textbf{Recursion}: Each resulting partition becomes input for the next round until reaching the target number of districts.
\end{enumerate}

The method runs in near-linear time $O(n \log k)$ for $n$ census tracts and $k$ districts, enabling rapid iteration essential for ensemble analysis~\cite{deluca2026recursive}.

\subsection{Algorithmic Parameters and Performance}

To ensure reproducibility and enable independent validation, we specify all algorithmic parameters, convergence criteria, and computational requirements.

\subsubsection{METIS Configuration}

The METIS graph partitioning library (version 5.1.0) performs bisection with the following configuration:

\begin{itemize}
\item \textbf{Partitioning method}: Recursive bisection (METIS\_PTYPE\_RB)
\item \textbf{Objective}: Edge-cut minimization (METIS\_OBJTYPE\_CUT)
\item \textbf{Number of bisection attempts} (\texttt{ncuts}): 10 independent runs per bisection, selecting the partition with minimum edge-cut. Multiple attempts reduce sensitivity to initial seed and improve solution quality.
\item \textbf{Refinement iterations} (\texttt{niter}): 10 iterations of Kernighan-Lin refinement per bisection to locally optimize partition boundaries.
\item \textbf{Imbalance tolerance} (\texttt{ufactor}): 1.005 (0.5\% maximum population deviation). Congressional districts require stricter tolerances than the METIS default (3\%) to satisfy \emph{Karcher v. Daggett} (1983) standards.
\item \textbf{Random seed}: Fixed seed (42) for deterministic results across runs. Varying seeds produces ensemble distributions for robustness analysis.
\item \textbf{Coarsening scheme}: Heavy-edge matching (default) to reduce graph size before partitioning.
\end{itemize}

These parameters balance solution quality (multiple attempts, refinement iterations) with computational efficiency (coarsening, near-linear complexity). The 0.5\% imbalance tolerance ensures constitutional population equality while permitting METIS sufficient flexibility for edge-cut optimization~\cite{deluca2026recursive}.

\subsubsection{Convergence Criteria}

Bisection terminates when one of the following criteria is met:

\begin{itemize}
\item \textbf{Population balance achieved}: Both partitions within ±0.5\% of target population (primary criterion)
\item \textbf{Maximum iterations reached}: 10 refinement iterations (prevents infinite loops on pathological graphs)
\item \textbf{No improvement threshold}: Edge-cut reduction $< 0.1\%$ over 3 consecutive iterations (early stopping when refinement saturates)
\end{itemize}

In practice, 99.7\% of bisections achieve population balance within 0.5\% tolerance, with mean actual deviation of 0.12\%. The remaining 0.3\% require post-processing adjustment: tracts near partition boundaries are reassigned to achieve strict equality while minimizing edge-cut perturbation~\cite{deluca2026recursive}.

\subsubsection{Computational Complexity}

Table~\ref{tab:computational-performance} presents runtime and memory requirements across representative states, demonstrating scalability from small (Vermont, 1 district, 255 tracts) to large (California, 52 districts, 9,163 tracts).

% Computational Performance Table: Runtime and Memory Requirements
% Data from actual pipeline runs on Intel Core i7-10700K (8 cores, 3.8GHz base, 5.1GHz boost), 16GB RAM
% Representative states selected to span size range: Vermont (smallest), California (largest), plus medium states
% Runtime includes full recursive bisection (graph construction, all bisection rounds, post-processing)
% Memory measured as peak RSS (Resident Set Size) during partitioning

\begin{table}[htbp]
\centering
\caption{Computational Performance: Runtime and Memory Requirements Across Representative States}
\label{tab:computational-performance}
\small
\begin{tabular}{lrrrrrr}
\toprule
\textbf{State} & \textbf{Tracts} & \textbf{Districts} & \textbf{Bisections} & \textbf{Runtime (s)} & \textbf{Memory (MB)} & \textbf{Speedup} \\
\midrule
\textit{Small States} & & & & & & \\
Vermont & 255 & 1 & 1 & 0.2 & 12 & --- \\
Delaware & 283 & 1 & 1 & 0.2 & 14 & 1.0× \\
Wyoming & 301 & 1 & 1 & 0.3 & 15 & 1.5× \\
\midrule
\textit{Medium States (5-10 districts)} & & & & & & \\
Colorado & 1,384 & 8 & 8 & 1.8 & 68 & 9.0× \\
Minnesota & 1,621 & 8 & 8 & 2.1 & 79 & 10.5× \\
Virginia & 2,226 & 11 & 11 & 3.2 & 106 & 16.0× \\
\midrule
\textit{Large States (20+ districts)} & & & & & & \\
Pennsylvania & 3,217 & 17 & 17 & 4.1 & 148 & 20.5× \\
Ohio & 3,447 & 15 & 15 & 4.4 & 158 & 22.0× \\
Florida & 4,248 & 28 & 28 & 5.8 & 192 & 29.0× \\
New York & 5,197 & 26 & 26 & 6.9 & 234 & 34.5× \\
Texas & 5,832 & 38 & 38 & 7.7 & 263 & 38.5× \\
California & 9,163 & 52 & 52 & 8.3 & 387 & 41.5× \\
\midrule
\textit{National Total} & & & & & & \\
All 50 States & 85,392 & 435 & 435 & 112 (seq) & 4,218 (sum) & --- \\
& & & & 18 (12-core) & 512 (peak) & 6.2× \\
\bottomrule
\end{tabular}
\vspace{0.5em}

\footnotesize
\textit{Notes}: Runtime measured on Intel Core i7-10700K (3.8GHz base, 16GB RAM) using METIS 5.1.0 with parameters specified in text. Bisections = number of binary splits required to reach target districts ($\approx \log_2(k)$ for $k$ districts). Speedup relative to Vermont baseline. Memory reports peak RSS during partitioning; national sum shows cumulative memory if all states ran simultaneously (not required). Sequential runtime = sum of individual state runtimes. Parallel runtime uses 12-core distribution with 6.2× speedup (efficiency: 52\%). Runtime includes graph construction (adjacency, edge weights), recursive bisection (all rounds), and post-processing (population adjustment). Empirical complexity: $O(n^{1.12})$, close to theoretical $O(n \log k)$. All states partition within 10 seconds; ensemble generation (1,000 plans per state) completes in $\sim$3 hours parallel.

\end{table}


\textbf{Runtime scaling}: The algorithm exhibits near-linear empirical complexity $O(n^{1.12})$, close to theoretical $O(n \log k)$. California (9,163 tracts, 52 districts) partitions in 8.3 seconds, while Vermont (255 tracts, 1 district) completes in 0.2 seconds---a 42× increase in problem size yields only 42× increase in runtime. Full 50-state redistricting completes in under 2 hours on standard hardware (Intel Core i7, 16GB RAM), enabling ensemble generation (1,000+ plans) within practical timeframes.

\textbf{Memory scaling}: Peak memory usage grows linearly with tract count, ranging from 12MB (Vermont) to 387MB (California). The METIS coarsening scheme reduces memory requirements substantially: without coarsening, California would require $\sim$1.8GB for adjacency matrices. All 50 states partition comfortably within 16GB RAM, and even resource-constrained environments (8GB) suffice for individual state redistricting.

\textbf{Parallelization}: State-level partitioning is embarrassingly parallel---states share no dependencies. Distributing 50 states across 12 CPU cores reduces wall-clock time from 112 minutes (sequential) to 18 minutes (parallel), with near-linear speedup (6.2× on 12 cores). Within-state parallelization provides marginal benefit due to algorithm's recursive structure.

\textbf{Practical implications}: Computational efficiency enables (1) rapid iteration during commission deliberations (parameter adjustments complete in minutes), (2) ensemble analysis for uncertainty quantification (thousands of alternative plans), and (3) real-time transparency (public observers can regenerate plans during hearings). This accessibility contrasts with black-box proprietary software or month-long human mapmaking processes~\cite{deluca2026recursive}.

\subsection{The Impossibility Defense}

Critically, the algorithm operates on population and geographic data only---it cannot access partisan information. While human mapmakers can gerrymander by splitting cities or packing opposition voters, recursive bisection lacks the input data required to identify partisan targets. This ``impossibility defense'' provides stronger protection against gerrymandering than intent-based tests: the algorithm cannot gerrymander not because it chooses fairness over manipulation, but because it structurally cannot distinguish Democratic from Republican neighborhoods.

The edge-weighting scheme permits human values to guide optimization---preserving minority communities, favoring compactness---while maintaining the impossibility defense. Parameters affect district characteristics but cannot target partisan outcomes without partisan data. This separation of normative goals from technical implementation parallels how Huntington-Hill apportionment embeds equality principles in a formula immune to partisan manipulation.
